% ===================================================================
% CHAPTER 5: ZEROTH-ORDER BASELINES
% ===================================================================
\chapter{Zeroth-Order Baselines}
\label{ch:zeroth}

% Ported from DETC2025-167604, "A Functional Indoor Testbed for Multirobot
% Adaptive Navigation in Vector Field Environments" (ref. 39), text at
% PhD Thesis/sources/idetc2025_testbed.txt. Its HSV encoding and NN
% calibration content is Chapter~\ref{ch:tools}'s. The unpublished 2024
% DETC2024-142746 manuscript is deliberately not used (author decision,
% 2026-09-29).

\section{Direction-Only Navigation}
\label{sec:ch5:direction_only}
% Ported from: idetc2025_testbed.txt Secs. 1, 2.5; opening paragraph new.

This chapter presents the zeroth-order control primitives of this
dissertation: primitives that command a cluster's heading from an
instantaneous, unweighted read of the team's vector-field measurements,
without fitting any spatial derivative of the field. Each robot returns a
single instantaneous vector reading of the local field, and the adaptive
navigation layer combines these readings into a commanded cluster
velocity, using the multilayer control architecture and cluster-space
formation control introduced in Chapters~\ref{ch:tools}
and~\ref{ch:estimation}. Robots are treated as omnidirectional; robot
orientation is not used, and the cluster's own orientation pose is not
commanded by any primitive in this chapter.

Two zeroth-order primitives from prior work \cite{38} are examined here:
the vector-sum technique, which uses the direction of the summed
readings, and the vector-to-scalar technique, which discards direction
and uses only the magnitude of each reading. Both were first validated
in simulation, with known limitations documented \cite{38}, and were
the first control primitives run on the Decabot testbed of
Chapter~\ref{ch:tools} \cite{39}
(Section~\ref{sec:ch5:hardware_trials}). Both are purely reactive: each
commands a velocity from the current instantaneous readings alone, with
no state carried between control cycles.

\section{Vector-Sum Primitive}
\label{sec:ch5:vector_sum}
% Ported from: idetc2025_testbed.txt Sec. 2.5.

The vector-sum technique combines the vectors sensed by each robot into
a resultant vector. The normalized direction of this resultant then
specifies the commanded cluster velocity $(\dot{x}_c,\dot{y}_c)$
\cite{39}. Because the readings are not normalized before summing, a
reading with larger magnitude carries proportionally more weight in the
resultant direction.

\section{Vector-to-Scalar Primitive}
\label{sec:ch5:vector_to_scalar}
% Ported from: idetc2025_testbed.txt Sec. 2.5, Eqs. 2-7.

In the vector-to-scalar technique, only the magnitude of each robot's
vector reading is used, reducing navigation to a scalar-field gradient
problem. Treating each robot's reading magnitude as a $z$-coordinate over
its $(x,y)$ position, the differences between robots' readings determine
the normal to the plane through the three points, which is then projected
to give the direction of steepest change. For robots $1$, $2$, $3$ at
positions $(x_i,y_i)$ with reading magnitudes $z_i$,
\begin{equation}
\vec{R}_{12} = \begin{pmatrix} x_2-x_1 \\ y_2-y_1 \\ z_2-z_1 \end{pmatrix}, \qquad
\vec{R}_{13} = \begin{pmatrix} x_3-x_1 \\ y_3-y_1 \\ z_3-z_1 \end{pmatrix},
\label{eq:ch5:R_edges}
\end{equation}
\begin{equation}
\vec{N} = -\vec{R}_{13} \times \vec{R}_{12},
\label{eq:ch5:normal}
\end{equation}
\begin{equation}
b_{\text{grad}} = \operatorname{ATAN2}(N_y, N_x) + \frac{\pi}{2},
\label{eq:ch5:bgrad}
\end{equation}
and the angle is converted to a commanded cluster velocity \cite{39},
\begin{equation}
\dot{x}_c = \cos(b_{\text{grad}}), \qquad
\dot{y}_c = \sin(b_{\text{grad}}).
\label{eq:ch5:vel_cmd}
\end{equation}
%% Corrected from source Eq. 7 (cos -> sin): as printed in
%% idetc2025_testbed.txt / DETC2025-167604, the pair reads
%% $\dot{x}_c = \cos(b_{\text{grad}})$ (Eq. 6, correct) and
%% $\dot{y}_c = \cos(b_{\text{grad}})$ (Eq. 7, printed with cos where sin
%% is needed for $(\dot{x}_c,\dot{y}_c)$ to be a unit vector at angle
%% $b_{\text{grad}}$). The affected equation is source Eq. 7, immediately
%% following the correctly-printed Eq. 6.

The gradient becomes ill-defined if the three robots are collinear.
Equation~\eqref{eq:ch5:bgrad} gives the direction of steepest ascent; the
technique navigates to a local extremum by moving in this direction
(source-seeking, or maximum-seeking) or its negation (sink-seeking, or
minimum-seeking, used for the hardware trials of
Section~\ref{sec:ch5:hardware_trials}), and orbits an extremum by moving
perpendicular to it.

\section{Hardware Trials on Fixed and Sinking Vortices}
\label{sec:ch5:hardware_trials}
% Ported from: idetc2025_testbed.txt Secs. 2.4-2.5 (trial cluster and
% primitive setup), 3.2 (Eqs. 8-13, field definitions), 3.5 (trial
% methodology), 4.2 (results); DETC2025-167604 \cite{39}. The
% Decabot testbed itself (HSV encoding, NN calibration, architecture) is
% Chapter~\ref{ch:tools}'s.

The vector-sum and vector-to-scalar primitives of
Sections~\ref{sec:ch5:vector_sum} and~\ref{sec:ch5:vector_to_scalar},
previously demonstrated only in simulation, were demonstrated on the
Decabot hardware testbed of Chapter~\ref{ch:tools} in \cite{39}, on
a fixed vortex and a sinking vortex field. The commanded formation was a
three-robot cluster with a constant shape policy, $p = q = 0.35$~m and
$\beta = 1.05$~rad, forming an equilateral triangle.

The first-order trials of Chapter~\ref{ch:first_order} used a smaller,
$0.33$~m equilateral formation.

Both fields are built from a fixed vortex centered at
$(\mathrm{center}_x,\mathrm{center}_y)$: with
$r = \sqrt{(x-\mathrm{center}_x)^2+(y-\mathrm{center}_y)^2}$ and
$\theta = \operatorname{atan2}(y-\mathrm{center}_y,\, x-\mathrm{center}_x)$,
\begin{equation}
u = r\sin\theta, \qquad v = -r\cos\theta.
\label{eq:ch5:fixed_vortex}
\end{equation}
The sinking vortex adds a radially inward term to this fixed vortex,
\begin{equation}
u \mathrel{+}= -\frac{c_s(x-\mathrm{center}_x)}{r^2+10^{-6}}, \qquad
v \mathrel{+}= -\frac{c_s(y-\mathrm{center}_y)}{r^2+10^{-6}},
\label{eq:ch5:sinking_vortex}
\end{equation}
a point sink of strength $c_s = 0.2$ whose inward speed decays as $1/r$, with
the $10^{-6}$ term added only to avoid division by zero at the center. After construction, both fields are rescaled so that vector
magnitude falls in $[0.3,1]$ for printing; magnitudes below 0.2 were
excluded from calibration, and below 0.3 did not produce a robot velocity
sufficient to overcome static friction.
%% Sign: DETC2025-167604 Eqs. 12-13 print the sink term with a plus sign
%% (outward), a typo; the hardware plotting script
%% trunk/robots_3/saddle_tests/plotter_avg_run_Sinking_Vortex.m uses the
%% inward form above, and the trials spiral inward.
%% c_s = 0.2 follows the testbed paper's own plotting package
%% (trunk/robots_3/Results_for TestBed_Paper/plotting_package/
%% plotter_avg_run_Sinking_Vortex.m); the published equations print 0.1.
%% Author decision 2026-09-29: use the code value.
%% This 1/r sink differs from the linear sinking vortex of
%% Chapter~\ref{ch:first_order} (Appendix~\ref{app:fields}).

Testing used a fixed 10~Hz OptiTrack sampling rate and 10 trial runs per
environment-primitive combination. The fixed vortex was tested with both
the vector-sum and vector-to-scalar (minimum-seeking) primitives; the
sinking vortex was tested only with vector-sum. Both fixed and randomized
starting positions were used for the vector-to-scalar trials on the fixed
vortex. Performance was evaluated by trajectory shape relative to the
expected path, steady-state error between final and target cluster
position, and formation-shape maintenance (Section~\ref{sec:ch5:formation_stats}).

On the fixed vortex, the vector-sum primitive produced a consistent
outward drift across all 10 runs, with the drift rate stable at all
angles, a behavior not predicted in simulation. Section~\ref{sec:ch5:why_not_location}
returns to this result. On the sinking vortex, the same primitive settled
the cluster to within approximately 0.1~m of the true center, with the
target location falling within the triangle connecting the three robots.
%% NOTE: the 0.2 m value comes from the text of DETC2025-167604 (third test
%% set, vector-to-scalar, fixed vortex, its Fig. 18), but that paper's Fig. 18
%% caption says "sinking vortex". Text and caption disagree in the published
%% paper; the text is followed here. Check against the Fig. 18 data.
The vector-to-scalar (minimum-seeking) primitive on the fixed vortex
settled the cluster within approximately 0.2~m of center, following an
approximately straight-line path with some deviation attributed to
printing noise and to variance in the saturation-to-magnitude estimate;
performance was essentially unchanged when starting positions were
randomized rather than fixed, since the velocity command is computed from
the local gradient of saturation rather than from saturation itself, and
near the center this gradient commands speeds below 0.05~m/s, too low
to overcome static friction as the robots were characterized at the
time (the floor was later remeasured at 0.025~m/s after recalibration,
Chapter~\ref{ch:tools}), regardless of where the cluster started. Neither orbiting nor saddle-point navigation, the subjects of
Chapter~\ref{ch:first_order}, was demonstrated on hardware at this stage;
both require the Jacobian-based estimation developed there.

\begin{figure}[h]
  \centering
  \begin{subfigure}[t]{0.32\textwidth}
    \centering
    \includegraphics[width=\textwidth]{ch5_fixed_vortex_vector_sum.png}
    \caption{Vector-sum, fixed vortex.}
  \end{subfigure}\hfill
  \begin{subfigure}[t]{0.32\textwidth}
    \centering
    \includegraphics[width=\textwidth]{ch5_sinking_vortex_vector_sum.png}
    \caption{Vector-sum, sinking vortex.}
  \end{subfigure}\hfill
  \begin{subfigure}[t]{0.32\textwidth}
    \centering
    \includegraphics[width=\textwidth]{ch5_fixed_vortex_vector_to_scalar.png}
    \caption{Vector-to-scalar, fixed vortex.}
  \end{subfigure}
  \caption{Average robot and cluster-center trajectories over 10 hardware
  runs per case, shaded to one standard deviation \cite{39}.}
  \label{fig:ch5:hardware_trial}
\end{figure}
% FIGURE SOURCE: trunk/robots_3/Results_for TestBed_Paper/ (Figs. 14, 15,
% 17 of DETC2025-167604): Fixed_vortex_Vector_Sum/average_trajectories.png,
% Sinking_Vortex_Vector_Sum/average_trajectories.png,
% Fixed_vortex_Vector_Sum/Fixed_Vortex_mag_Descent_consistent_Start/average_trajectories.png

\begin{figure}[h]
  \centering
  \begin{subfigure}[t]{0.48\textwidth}
    \centering
    \includegraphics[width=\textwidth]{ch5_sinking_vortex_distance.png}
    \caption{Vector-sum, sinking vortex.}
  \end{subfigure}\hfill
  \begin{subfigure}[t]{0.48\textwidth}
    \centering
    \includegraphics[width=\textwidth]{ch5_vector_to_scalar_distance.png}
    \caption{Vector-to-scalar, fixed vortex.}
  \end{subfigure}
  \caption{Average distance of the cluster center from the field center
  over time, 10 hardware runs per case \cite{39}.}
  \label{fig:ch5:hardware_distance}
\end{figure}
% FIGURE SOURCE: Figs. 16 and 18 of DETC2025-167604; average_distance_vs_time.png
% in the same two folders.

\section{Formation Maintenance Statistics}
\label{sec:ch5:formation_stats}
% Ported from: idetc2025_testbed.txt Sec. 4.3, Table 2 \cite{39}.

Formation shape was held close to its commanded values throughout the
hardware trials of Section~\ref{sec:ch5:hardware_trials}
(Table~\ref{tab:ch5:formation}) \cite{39}, with deviations on the
order of 1~cm in $p$ and $q$, consistent with the resolution limits of
the OptiTrack motion capture system itself (Chapter~\ref{ch:tools}).

\begin{table}[h]
  \centering
  \caption{Cluster formation-holding accuracy across control primitive
  trials, commanded $p=q=0.35$~m, $\beta=1.05$~rad, 10 trial runs per
  column \cite{39}.}
  \label{tab:ch5:formation}
  \begin{tabular}{lccc}
    \toprule
     & $p$ (m) & $q$ (m) & $\beta$ (rad) \\
    \midrule
    Fixed vortex, vector-sum & $0.36 \pm 0.028$ & $0.36 \pm 0.032$ & $1.03 \pm 0.122$ \\
    Fixed vortex, vector-to-scalar & $0.35 \pm 0.015$ & $0.36 \pm 0.009$ & $1.02 \pm 0.045$ \\
    Sinking vortex, vector-sum & $0.35 \pm 0.019$ & $0.35 \pm 0.015$ & $1.02 \pm 0.056$ \\
    \bottomrule
  \end{tabular}
\end{table}

\section{Why Direction Is Not Location}
\label{sec:ch5:why_not_location}
% New text.

Every primitive in this chapter commands a heading. None of them
estimates a position: nothing in Sections~\ref{sec:ch5:vector_sum}
or~\ref{sec:ch5:vector_to_scalar} produces a coordinate for the feature
being sought, only a direction to move in on the current control cycle.
On a purely rotational field component, this is a structural limitation
rather than a tuning problem. A vortex's field vectors are everywhere
tangential to circles about its center, so a primitive that follows the
locally sensed direction, as vector-sum does, tracks that tangential direction with no term correcting
for radial distance from the center. The fixed-vortex hardware result of
Section~\ref{sec:ch5:hardware_trials}, a consistent outward drift for the
vector-sum primitive across all 10 runs and stable at all angles, is this
failure mode observed directly, and on hardware it was worse than
simulation predicted: \cite{39} attributes the drift to the robots
always moving tangent to the circle while momentum effects, not captured
in the simulation model, carry them further outward on each cycle.
Orbiting at a fixed radius is reported there as an open problem the
testbed is positioned to address, not one these primitives solve.

The same qualitative failure appears in the critical-points paper's own
comparison of primitives (Figure~\ref{fig:ch6:primitive_comparison}): both the
vector-sum and vector-to-scalar primitives exhibit outward spiral drift
in a vortex field when compared against a Jacobian-based orbital
controller, because neither has a mechanism to detect or correct radial
error with respect to a specific point. In a saddle field the comparison
is more decisive still: the field vectors do not all point toward the
critical point, so the vector-sum primitive cannot produce an orbital
trajectory at all.


This is why direction-only control cannot certify arrival at a point:
arrival is a statement about position, and a heading command carries no
position information to certify it with. The vector-to-scalar primitive
comes closest, since it converts the problem to scalar-field gradient
ascent or descent and can locate an extremum of the magnitude field, but
it still has no notion of the field's vector structure away from that
extremum, and so cannot classify what kind of critical point it has
found or navigate to a saddle. Recovering that structure requires
estimating the field's local Jacobian, which is what makes a critical
point's location, rather than merely its uphill direction, computable
(Chapter~\ref{ch:first_order}). The same three-robot cluster suffices
for that estimate; the change is that the readings are fitted to an
affine model instead of summed.

\section{Summary}
\label{sec:ch5:summary}
% New text.

This chapter presented two zeroth-order control primitives: vector-sum,
built from the direction of the summed readings, and vector-to-scalar,
built from their magnitudes alone. On 10 hardware trial runs per
environment-primitive combination \cite{39}, vector-sum settled near the
center of a sinking vortex and vector-to-scalar settled near the center
of a fixed vortex, with commanded formation shape held to within about
1~cm. What they cannot do is certify that the cluster has arrived
anywhere: with no position estimate, a rotational field component drives
a steady outward drift that these primitives have no mechanism to
correct, worse on hardware than simulation predicted.
Chapter~\ref{ch:first_order} takes the next step in the ladder of local
polynomial fit order, estimating the field's Jacobian from the same
three robots to recover a critical point's location and type rather than
only its local uphill direction.
